%%=============================================================================
%% Voorwoord
%%=============================================================================

\chapter*{\IfLanguageName{dutch}{Woord vooraf}{Preface}}%
\label{ch:voorwoord}

%% TODO:
%% Het voorwoord is het enige deel van de bachelorproef waar je vanuit je
%% eigen standpunt (``ik-vorm'') mag schrijven. Je kan hier bv. motiveren
%% waarom jij het onderwerp wil bespreken.
%% Vergeet ook niet te bedanken wie je geholpen/gesteund/... heeft

Voor u ligt de bachelorproef "Optimalisatie van personeelsplanning in lokale bakkerijen met behulp van operationele data en machine-learning". Dit werk vormt de afsluiting van mijn opleiding Toegepaste Informatica aan de Hogeschool Gent.

De keuze voor dit onderwerp kwam voort uit een persoonlijke motivatie om de digitale kloof voor kleine ondernemingen te verkleinen. Bakkerij Bossu in Nazareth vormde hiervoor de ideale casus. Ik heb van dichtbij gezien hoe uitdagend de dagelijkse personeelsplanning kan zijn in een sector die zo sterk onderhevig is aan externe factoren zoals feestdagen en seizoensinvloeden. Met dit onderzoek hoop ik een brug te slaan tussen complexe data-analyse en de praktische realiteit van de ambachtelijke bakker.

Een woord van dank is hier op zijn plaats. In het bijzonder wil ik mijn co-promotor, Dhr. S. Bossu, bedanken voor de waardevolle begeleiding en het gestelde vertrouwen in dit onderzoek. Zijn data en inzichten hebben dit project naar een hoger niveau getild.

\vspace{1cm}

\noindent Gaétan Dessomviele